% -----------------------------------------------------------
\section{Sin}
% -----------------------------------------------------------
	\label{SIN}
	
% ----------------------
\subsection{See also}
% ----------------------

	\hyperref[ACCOUNTABILITY]{Accountability}, 

% % ----------------------
% \subsection{1828 American Dictionary of the English Language} % *1828
% % ----------------------
	% \label{SIN-1828}

	% \begin{compactitem}
		% \item
	% \end{compactitem}

% % ----------------------
% \subsection{Oxford English Dictionary} % *OED
% % ----------------------
	% \label{SIN-OED}

	% \begin{compactitem}
		% \item[]
	% \end{compactitem}

% ----------------------
\subsection{Scriptural Usage} % *SCRIPTURES
% ----------------------
	\label{SIN-SCRIPTURES}

	% % -----
	% \subsubsection{Old Testament} % *OT
	% % -----
		% \label{SIN-OT}
	
		% % --
		% \paragraph{KJV}
		% % --
			% \label{SIN-OT-KJV}
	
			% \begin{tabular}{ll}
				% Total occurrences in the OT & 
			% \end{tabular}
			
			% \begin{compactdesc}
				% \item[]
			% \end{compactdesc}
			
		% % --
		% \paragraph{Hebrew Bible} % *HEBREW
		% % --
			% \label{SIN-HEBREW}
		
			% %
			% \subparagraph{\texorpdfstring{\he{sample word}}{}}
			% %
				% \label{PAGA} % whatever is in step bible without periods
			
				% \begin{compactdesc}
					% \item[HALOT , PDF ] \textbf{}
						% \begin{compactitem}
							% \item[1]
						% \end{compactitem}
					% \item[BDB , PDF ] \textbf{}
						% \begin{compactitem}
							% \item[1]
						% \end{compactitem}
					% \item[DCH PDF ] \textbf{}
						% \begin{compactitem}
							% \item[1]
						% \end{compactitem}
				% \end{compactdesc}
				
				% \begin{compactdesc}
					% \item[Some OT uses] \textbf{}
						% \begin{compactdesc}
							% \item[]
						% \end{compactdesc}
				% \end{compactdesc}
			
		% % --
		% \paragraph{Septuagint}
		% % --
			% \label{SIN-LXX}
		
			% %
			% \subparagraph{\texorpdfstring{\gr{sample word}}{}}
			% %
		
	% -----
	\subsubsection{New Testament} % *NT
	% -----
		\label{SIN-NT}
	
		% --
		\paragraph{KJV}
		% --
			\label{SIN-NT-KJV}
			
	
			\begin{tabular}{ll}
				Total occurrences in the NT & 227 % "sin" + "sinful" + "sinning" + "sinned" + "sinneth" + "sinnest" + "sins"
			\end{tabular}
			
			\begin{compactdesc}
				\item[{\hyperref[MATTHEW-9-5]{Matthew 9:5--6}}]
				\item[{\hyperref[MATTHEW-26-28]{Matthew 26:28}}]
				\item[{\hyperref[MARK-2-10]{Mark 2:10}}]
				\item[{\hyperref[LUKE-5-24]{Luke 5:24}}]
				\item[{\hyperref[JOHN-8-34]{John 8:34}}]
				\item[{\hyperref[JOHN-15-22]{John 15:22--24}}] \phantom{.}
					\begin{compactitem}
						\item Note the word \gr{πρόφασιν} here, interesting and important word.
					\end{compactitem}
				\item[{\hyperref[JOHN-20-23]{John 20:23}}]
				\item[{\hyperref[ROMANS-2-12]{Romans 2:12}}]
				\item[{\hyperref[ROMANS-3-19]{Romans 3:19--31}}] 
				\item[{\hyperref[ROMANS-4-7]{Romans 4:7--8}}] 
				\item[{\hyperref[ROMANS-5-12]{Romans 5:12--17}}] 
				\item[{\hyperref[ROMANS-5-20]{Romans 5:20--21}}] 
				\item[{\hyperref[ROMANS-6]{Romans 6--8}}] \phantom{.}
					\begin{compactitem}
						\item So many different things in these chapters.
					\end{compactitem}
				\item[{\hyperref[ROMANS-14-23]{Romans 14:23}}]
				\item[{\hyperref[1CORINTHIANS-15-56]{1 Corinthians 15:56}}]
				\item[{\hyperref[HEBREWS-10-17]{Hebrews 10:17}}]
				\item[{\hyperref[JAMES-1-15]{James 1:15}}]
				\item[{\hyperref[JAMES-4-17]{James 4:17}}]
				\item[{\hyperref[1JOHN-1-6]{1 John 1:6--2:2}}]
				\item[{\hyperref[1JOHN-3-4]{1 John 3:4--10}}]
				\item[{\hyperref[1JOHN-5-17]{1 John 5:17}}]
			\end{compactdesc}
			
		% --
		\paragraph{Greek New Testament} % *GREEK
		% --
			\label{SIN-GREEK}
			
			%
			\subparagraph{\texorpdfstring{\gr{ἁμαρτάνω}}{}}
			%
				\label{HAMARTANO}
			
				\begin{compactdesc}
					\item[BDAG 43b] \textbf{}
						\begin{compactitem}
							\item[1] to commit a wrong, \textit{to sin}
							\item[1.a] absolute
							\item[1.b] with fuller indication of that in which the mistake or moral failure consists, by means of a supplementary participle
							\item[1.c] with indication of the manner of sinning
							\item[1.d] with indication of the one against whom the sin is committed
							\item[1.e] with indication of the result (cf. BDAG \#1b for \gr{ἁμαρτία})
						\end{compactitem}
					\item[LSJ 77a, PDF 150] \textbf{}
						\begin{compactitem}
							\item[I.1] \textit{miss the mark}, especially of a spear thrown
							\item[I.2] generally, \textit{fail of one's purpose, go wrong}
							\item[I.3] fail of having, be deprived of
							\item[I.4] rarely, \textit{fail to do, neglect}
							\item[II.1] do wrong, err, sin
							\item[II.2] passive, \textit{to be frustrated} (of plans, etc.)
							\item[II.3] participle as adjective, \textit{wrong, mistaken}
						\end{compactitem}
				\end{compactdesc}
				
				\begin{compactdesc}
					\item[Some NT uses] \textbf{}
						\begin{compactdesc}
							\item[]
						\end{compactdesc}
				\end{compactdesc}
		
			%
			\subparagraph{\texorpdfstring{\gr{ἁμαρτία}}{}}
			%
				\label{HAMARTIA}
				
				See also \hyperref[PARABASIS]{\grl{παράβασις}}, \hyperref[ANOMIA]{\grl{ἀνομία}}, \hyperref[ADIKIA]{\grl{ἀδικία}}, 
			
				\begin{compactdesc}
					\item[BDAG 44a] \textbf{}
						\begin{compactitem}
							\item[1] a departure from either human or divine standards of uprightness
							\item[1.a] \textit{sin}
							\item[1.b] special sins (of apostasy, sin that leads to death, etc.)
							\item[2] a state of being sinful, \textit{sinfulness}
							\item[3] A destructive evil power, \textit{sin}
							\item[3.a] Paul thinks of sin almost in personal terms with other divinities of the nether world, as a ruling power that invades the world.
							\item[3.b] In Hebrews (as in the OT) sin appears as the power that deceives humanity and leads it to destruction, whose influence and activity can be ended only by sacrifices.
						\end{compactitem}
					\item[LSJ 77b, PDF 150] \textbf{}
						\begin{compactitem}
							\item[1] a failure, fault; error of judgment
							\item[2] in philosophy and religion, \textit{guilt, sin}
						\end{compactitem}
				\end{compactdesc}
				
				\begin{compactdesc}
					\item[Some NT uses] \textbf{}
						\begin{compactdesc}
							\item[]
						\end{compactdesc}
				\end{compactdesc}
		
	% -----
	\subsubsection{Book of Mormon} % *BOM
	% -----
		\label{SIN-BOM}
	
		\begin{tabular}{ll}
			Total occurrences in the BOM & 235 % "sin" + "sinful" + "sinning" + "sinned" + "sinneth" + "sinnest" + "sins"
		\end{tabular}
		
		\begin{compactdesc}
			\item[2 Nephi 2:13] And if ye shall say there is no law, ye shall also say there is no sin. And if ye shall say there is no sin, ye shall also say there is no righteousness. And if there be no righteousness, there be no happiness. And if there be no righteousness nor happiness, there be no punishment nor misery. And if these things are not, there is no God. And if there is no God, we are not, neither the earth. For there could have been no creation of things, neither to act nor to be acted upon; wherefore all things must have vanished away.
			\item[2 Nephi 2:23] and they would have had no children. Wherefore they would have remained in a state of innocence, having no joy, for they knew no misery, doing no good, for they knew no sin.
			\item[Mosiah 3:11] For behold, and also his blood atoneth for the sins of those who have fallen by the transgression of Adam who hath died not knowing the will of God concerning them or who have ignorantly sinned.
			\item[Mosiah 3:16] And even if it were possible that little children could sin, they could not be saved. But I say unto you: They are blessed; for behold, as in Adam or by nature they fall, even so the blood of Christ atoneth for their sins.
			\item[Mosiah 4:29] And finally, I cannot tell you all the things whereby ye may commit sin; for there are divers ways and means, even so many that I cannot number them.
			\item[Alma 28:13] And thus we see how great the unequality of man is because of sin and transgression and the power of the devil, which comes by the cunning plans which he hath devised to ensnare the hearts of men.
			\item[{\hyperref[ALMA-29-3]{Alma 29:3}}]
			\item[Alma 41:10] Do not suppose because that it hath been spoken concerning restoration that ye shall be restored from sin to happiness. Behold, I say unto you: Wickedness never was happiness.
			\item[Alma 42:16--24] Now repentance could not come unto men except there were a punishment, which also was as eternal as the life of the soul, should be affixed opposite to the plan of happiness, which was as eternal also as the life of the soul. 17 Now how could a man repent except he should sin? How could he sin if there was no law? How could there be a law save there was a punishment? 18 Now there was a punishment affixed and a just law given, which brought remorse of conscience unto man. 19 Now if there was no law given if a man murdered he should die, would he be afraid he should die if he should murder? 20 And also if there was no law given against sin, men would not be afraid to sin. 21 And if there was no law given if men sinned, what could justice do, or mercy either? For they would have no claim upon the creature. 22 But there is a law given and a punishment affixed and repentance granted, which repentance mercy claimeth. Otherwise justice claimeth the creature and executeth the law, and the law inflicteth the punishment. If not so, the works of justice would be destroyed; and God would cease to be God. 23 But God ceaseth not to be God; and mercy claimeth the penitent, and mercy cometh because of the atonement, and the atonement bringeth to pass the resurrection of the dead, and the resurrection of the dead bringeth back men into the presence of God. And thus they are restored into his presence, to be judged according to their works, according to the law and justice. 24 For behold, justice exerciseth all his demands; and also mercy claimeth all which is her own. And thus none but the truly penitent are saved.
			\item[{\hyperref[ALMA-45-16]{Alma 45:16}}]
			\item[Helaman 5:10--11] And remember also the words which Amulek spake unto Zeezrom in the city of Ammonihah, for he said unto him that the Lord surely should come to redeem his people, but that he should not come to redeem them in their sins but to redeem them from their sins. 11 And he hath power given unto him from the Father to redeem them from their sins because of repentance. Therefore he hath sent his angels to declare the tidings of the conditions of repentance, which bringeth unto the power of the Redeemer, unto the salvation of their souls.
			\item[3 Nephi 6:18] Now they did not sin ignorantly; for they knew the will of God concerning them, for it had been taught unto them. Therefore they did willfully rebel against God.
			\item[Mormon 2:13] But behold, this my joy was vain; for their sorrowing was not unto repentance because of the goodness of God, but it was rather the sorrowing of the damned because the Lord would not always suffer them to take happiness in sin.
			\item[Moroni 7:12] Wherefore all things which are good cometh of God; and that which is evil cometh of the devil. For the devil is an enemy unto God and fighteth against him continually and inviteth and enticeth to sin and to do that which is evil continually.
			\item[Moroni 8:11] And their little children need no repentance, neither baptism. Behold, baptism is unto repentance, to the fulfilling the commandments unto the remission of sins.
			\item[Moroni 8:25--26] And the firstfruits of repentance is baptism. And baptism cometh by faith unto the fulfilling the commandments; and the fulfilling the commandments bringeth remission of sins; 26 and the remission of sins bringeth meekness and lowliness of heart. And because of meekness and lowliness of heart cometh the visitation of the Holy Ghost, which Comforter filleth with hope and perfect love, which love endureth by diligence unto prayer until the end shall come, when all the saints shall dwell with God.
		\end{compactdesc}
		
	% -----
	\subsubsection{Doctrine and Covenants} % *DC
	% -----
		\label{SIN-DC}
	
		\begin{tabular}{ll}
			Total occurrences in the DC & 146 % "sin" + "sinful" + "sinning" + "sinned" + "sinneth" + "sinnest" + "sins"
		\end{tabular}
		
		\begin{compactdesc}
			\item[{\hyperref[DC1-31]{DC 1:31}}] For I the Lord cannot look upon sin with the least degree of allowance;
			\item[{\hyperref[DC29-46]{DC 29:46--47}}]
			\item[{\hyperref[DC45-4]{DC 45:4}}]
			\item[{\hyperref[DC49-8]{DC 49:8}}] 
			\item[{\hyperref[DC58-15]{DC 58:15}}] 
			\item[{\hyperref[DC58-42]{DC 58:42--43}}] 
			\item[{\hyperref[DC59-13]{DC 59:13--15}}]
			\item[{\hyperref[DC64-7]{DC 64:7--9}}]
			\item[{\hyperref[DC68-25]{DC 68:25}}]
			\item[{\hyperref[DC82-3]{DC 82:3--4}}]
			\item[{\hyperref[DC82-7]{DC 82:7}}]
			\item[{\hyperref[DC84-48]{DC 84:48--53}}] 
			\item[{\hyperref[DC88-32]{DC 88:32--40}}]
			\item[{\hyperref[DC93-1]{DC 93:1}}] 
			\item[{\hyperref[DC95-1]{DC 95:1--3}}]
				\begin{compactitem}
					\item Especially verse 3---the sin here is not following an important commandment that they were given.
				\end{compactitem}
			\item[{\hyperref[DC121-34]{DC 121:34--46}}] 
			\item[{\hyperref[DC138-57]{DC 138:57}}] 
		\end{compactdesc}
		
	% -----
	\subsubsection{Pearl of Great Price} % *PGP
	% -----
		\label{SIN-PGP}
	
		\begin{tabular}{ll}
			Total occurrences in the PGP & 17 % "sin" + "sinful" + "sinning" + "sinned" + "sinneth" + "sinnest" + "sins"
		\end{tabular}
		
		\begin{compactdesc}
			\item[Moses 6:54--56] Hence came the saying abroad among the people, that the Son of God hath atoned for original guilt, wherein the sins of the parents cannot be answered upon the heads of the children, for they are whole from the foundation of the world. 55 And the Lord spake unto Adam, saying: Inasmuch as thy children are conceived in sin, even so when they begin to grow up, sin conceiveth in their hearts, and they taste the bitter, that they may know to prize the good. 56 And it is given unto them to know good from evil; wherefore they are agents unto themselves, and I have given unto you another law and commandment. 
		\end{compactdesc}
		
% ----------------------
\subsection{Key Phrases} % *KEYPHRASES *PHRASES
% ----------------------
	\label{SIN-PHRASES}

	\begin{compactdesc}
		\item[under sin] \textbf{5} | \hyperref[ROMANS-3-9]{Romans 3:9}; \hyperref[ROMANS-7-14]{Romans 7:14}; \hyperref[GALATIANS-3-22]{Galatians 3:22}; \hyperref[DC49-8]{DC 49:8}; \hyperref[DC84-53]{DC 84:53}
			\begin{compactdesc}
				\item[Bible anchor verses] \phantom{.}
					\begin{compactdesc}
						\item[Romans 3:9] \gr{Τί οὖν; προεχόμεθα; οὐ πάντως, προῃτιασάμεθα γὰρ Ἰουδαίους τε καὶ Ἕλληνας πάντας ὑφ᾽ ἁμαρτίαν εἶναι,}
						\item[Romans 7:14] \gr{Οἴδαμεν γὰρ ὅτι ὁ νόμος πνευματικός ἐστιν· ἐγὼ δὲ σάρκινός εἰμι, πεπραμένος ὑπὸ τὴν ἁμαρτίαν.}
						\item[Galatians 3:22] \gr{ἀλλὰ συνέκλεισεν ἡ γραφὴ τὰ πάντα ὑπὸ ἁμαρτίαν ἵνα ἡ ἐπαγγελία ἐκ πίστεως Ἰησοῦ Χριστοῦ δοθῇ τοῖς πιστεύουσιν.}
					\end{compactdesc}
				% \item[Commentary] ---
			\end{compactdesc}
	\end{compactdesc}
	
% % ----------------------
% \subsection{Bible Dictionaries} % *DICTIONARIES
% % ----------------------
	% \label{SIN-BD}

% % ----------------------
% \subsection{Restoration Usage} % *RESTORATION
% % ----------------------
	% \label{SIN-RESTORATION}

	% % -----
	% \subsubsection{Joseph Smith} % *JS
	% % -----
		% \label{SIN-JS}
	
		% \begin{compactdesc}
			% \item[DATE/SOURCE]
		% \end{compactdesc}
	
	% % -----
	% \subsubsection{Brigham Young} % *BY
	% % -----
		% \label{SIN-BY}
	
		% \begin{compactdesc}
			% \item[DATE/SOURCE]
		% \end{compactdesc}
	
% ----------------------
\subsection{Commentary and Studies} % *COMMENTARY *STUDIES
% ----------------------
	\label{SIN-COMMENTARY}
	
	% -----
	\subsubsection{Key Points}
	% -----
		\label{SIN-KEYS}
	
		\begin{compactitem}
			\item Definitions/understandings of sin
				\begin{compactitem}
					\item \hyperref[JAMES-4-17]{James 4:17}
						\begin{compactitem}
							\item My translation: ``So, for a person who is able to do good, and doesn't do it---for that person, it is sin.''
							\item ``Good'' has a very expansive meaning here, referring to anything that could contribute to salvation.
							\item See \hyperref[SIN-james-4-17]{this study} below.
						\end{compactitem}
					\item \hyperref[1JOHN-3-4]{1 John 3:4}
						\begin{compactitem}
							\item My translation: ``Everyone who is practicing sin is also practicing lawlessness; and sin is lawlessness.''
							\item See \hyperref[SIN-1john-3-4]{this study} below.
						\end{compactitem}
					\item \hyperref[1JOHN-5-17]{1 John 5:17}
						\begin{compactitem}
							\item The key part of the verse is that all \gr{ἀδικία} is sin.
							\item \gr{ἀδικία} relates to laws, righteousness, and \gr{ἀνομία}; see \hyperref[SIN-1john-5-17]{this study} below.
						\end{compactitem}
				\end{compactitem}
		\end{compactitem}
	
	% -----
	\subsubsection{What is sin, and how does it relate to righteousness?}
	% -----
		\label{SIN-definition}
		
		The scriptures have a lot to say about sin,
			and the term appears in many discussions of the plan of salvation.
		It's related to many other important things in the gospel as well,
			including righteousness,
				laws,
					and repentance.
		Here I'll use the scriptures to offer a definition of sin as well as an explanation of how it relates to other aspects of the gospel.
		The definition and my explanation are ways of understanding sin;
			they are not the only ways,
				though I've tried to include as much scriptural evidence as possible.
		
		% --
		\paragraph{What is sin?}
		% --
		
			The restoration scriptures don't give a clear definition of ``sin.''
			We instead get many examples of individual sins,
				along with discussions of the consequences of sin.
			Examples of something and its consequences are helpful for grasping it,
				but it would also be nice to have a more general statement of how to understand sin,
					independent of particular examples.
			
			Several verses in the New Testament discuss sin in this more general way.
			The three most important are \hyperref[JAMES-4-17]{James 4:17},
				\hyperref[1JOHN-3-4]{1 John 3:4},
					and \hyperref[1JOHN-5-17]{1 John 5:17}.
			I'll look at these verses one at a time,
				and then we'll see what they tell us together.
				
			%
			\subparagraph{\hyperref[JAMES-4-17]{James 4:17}}
			%
				\label{SIN-james-4-17}
			
				The first verse is James 4:17.
				Here is the verse in the KJV,
					along with the Greek text.
					
					\begin{description}
						\item[KJV] Therefore to him that knoweth to do good, and doeth it not, to him it is sin.
						\item[SBLG] \gr{εἰδότι οὖν καλὸν ποιεῖν καὶ μὴ ποιοῦντι, ἁμαρτία αὐτῷ ἐστιν.}
					\end{description}
				
				The first word in Greek,
					\gr{εἰδότι},
						is a participle of the verb \gr{οἶδα},
							``to know.''
				The participle is in dative case and so becomes the recipient or owner of the \gr{ἁμαρτία},
					the sin,
						in the second part of the verse;
							the pronoun \gr{αὐτῷ} at the end is also in dative case and its antecedent is \gr{εἰδότι}.
				The second participle in the verse,
					\gr{ποιοῦντι},
						also makes clear that the person who knows is the same person as the one who isn't acting on that knowledge.
						
				The KJV is a fine translation but obscures some aspects of the verse.
				Here are some other translations that will be helpful later:
				
					\begin{description}
						\item[Anchor Bible] Therefore it counts as a sin for the person who understands the proper thing to do and yet does not do it.
						\item[Syriac-English] and anyone who does not carry out the things that they know to be good, is really committing a sin.
						\item[NASB] Therefore, to one who knows the right thing to do and does not do it, to him it is sin.
						\item[NIV] If anyone, then, knows the good they ought to do and doesn't do it, it is sin for them.
						\item[NRSV] Anyone, then, who knows the right thing to do and fails to do it, commits sin.
					\end{description}
					
				An interpretation of the Greek which we can rule out is one which takes the ``knowing'' here as having some sort of normative force,
					where we read the verse as saying that someone \textit{knows to do} something,
						or \textit{knows that they should} do something.
				Sometimes,
					in English,
						we use ``know'' this way.
				I can say that a friend ``knows to arrive at the restaurant at 6:30,''
					by which I mean that the person knows they \textit{should} arrive at that time,
						or that we are \textit{expecting} them to arrive at that time.
				This use of ``know'' has normative force,
					but that is not what the word means in this verse.
				There is no support for reading \gr{οἶδα} this way,
					so we ought to reject this interpretation,
						which is based on an unusual way of using the English word to begin with.
						
				Some of the translations listed above do give normative force to the verse,
					but the normativity comes not from the participle \gr{εἰδότι} but instead from \gr{καλὸν},
						the ``good thing'' or ``right thing.''
				\gr{εἰδότι} completes its meaning with the infinitive \gr{ποιεῖν};
					the object of \gr{ποιεῖν} is \gr{καλὸν}.
				
				Forms of \gr{οἶδα},
					``to know,''
						can take an infinitive,
							as \gr{εἰδότι} does here.
				The proper sense of the participle here,
					then,
						is the third meaning in BDAG,
							or ``to know/understand how, \textit{can, be able}'' (see these definitions \hyperref[OIDA]{here}).
				BDAG notes that this use of \gr{οἶδα} takes an infinitive following it,
					which is what we get here,
						and it also lists this verse in James an example of this sense of the word (this meaning corresponds to \#10 in CGL and \#I.2 in LSJ).
				So our interpretations of the verse should focus on this sense of the participle \gr{εἰδότι}.
				
				Here is a way of rendering the verse to bring out that meaning:
				
					\begin{quote}
						So, for a person who knows how to do good, and doesn't do it---for that person, it is sin.
					\end{quote}
					
				\noindent In this translation,
					sin is when someone knows how to do something good,
						but for some reason doesn't do it. 
				Maybe they chose a bad thing instead,
					or just did nothing.
				Either way there was something good that they knew how to do and didn't do it,
					so it was sin.
					
				Knowing how to do something is closely related to being \textit{able} to do something---%
					that is,
						usually at least,
							part of being able to do something is knowing how to do it (another part might be, there's no one stopping you from doing it).
				This is why one meaning of \gr{οἶδα} grows out of the ``knowing how'' sense to involve an ability or possession of a skill,
					thereby making someone able to do something (the same meanings listed above in the lexicons include this sense of being able).
				Here is another possible translation along those lines:
				
					\begin{quote}
						So, for a person who is able to do good, and doesn't do it---for that person, it is sin.
					\end{quote}
					
				\noindent The only difference between this translation and the previous one is that we're talking in terms of ``being able'' instead of ``knowing how.''
				As I said,
					those aren't the same thing,
						but they are similar.
				To be able to do good,
					you must not only know how,
						but also be in a position where doing good is possible (you aren't tied up in a chair,
							for example).
				To know how to do good,
					all you need is the knowledge of how to do good.
				You can meet that requirement when you're tied up in a chair,
					but I don't think it makes sense to say that such a person has sinned when they don't do something good,
						or that their lack of action is a sin.
				They aren't able to act and so couldn't have sinned.
				Putting this verse in terms of being able to act lets us get \gr{οἶδα}'s sense of ``knowing how,''
					while also making sure that we don't charge sin to people who aren't actually able to do something.
				On this second translation of the verse,
					sin is when a person (1) knows \textit{how} to do good,
						(2) is \textit{able} to do good,
							and (3) still doesn't actually do any good.
							
				The last thing to discuss with this verse is \gr{καλὸν},
					or what is meant by ``good'' or ``a good thing.''
				In Greek there are two words which we usually translate as ``good''---%
					\hyperref[KALOS]{\grl{καλός}} and \hyperref[AGATHOS]{\grl{ἀγαθός}}.
				While they both can mean ``good'' in the sense of ``useful, beneficial'' and ``good'' in the sense of ``morally good,''
					\gr{καλός} can also mean ``beautiful, good-looking'' or ``noble.''
				BDAG meaning \#2.b,
					under which it lists James 4:17,
						also adds ``praiseworthy'' and ``contributing to salvation'' as possible meanings.
				In this understanding of sin,
					then,
						we have to keep the expansive meaning of \gr{καλὸν} in mind---%
							it refers to any good,
								excellent,
									or contributing-to-salvation thing we could do.
				We find this expansive meaning in some restoration scripture,
					such as \hyperref[MORONI-7-13]{Moroni 7:13} and \hyperref[MORONI-10-6]{Moroni 10:6}.
					
				Here is the translation of James 4:17 one more time:
				
					\begin{quote}
						So, for a person who is able to do good, and doesn't do it---for that person, it is sin.
					\end{quote}
			
			%
			\subparagraph{\hyperref[1JOHN-3-4]{1 John 3:4}}
			%
				\label{SIN-1john-3-4}
				
				The next verse to consider is 1 John 3:4.
				Here is the KJV translation,
					along with the Greek:
					
					\begin{description}
						\item[KJV] Whosoever committeth sin transgresseth also the law: for sin is the transgression of the law.
						\item[SBLG] \gr{Πᾶς ὁ ποιῶν τὴν ἁμαρτίαν καὶ τὴν ἀνομίαν ποιεῖ, καὶ ἡ ἁμαρτία ἐστὶν ἡ ἀνομία.}
					\end{description}
					
				\noindent The KJV translation is a little more helpful for this verse,
					but here are some alternatives for comparison:
					
					\begin{description}
						\item[Anchor Bible] Everyone who acts sinfully is really doing iniquity,
									for sin is the Iniquity.
						\item[Syriac-English] But whoever does what is sinful commits iniquity, for every sin is iniquity.
						\item[NASB] Everyone who practices sin also practices lawlessness; and sin is lawlessness.
						\item[NIV] Everyone who sins breaks the law; in fact, sin is lawlessness.
						\item[NRSV] Everyone who commits sin is guilty of lawlessness; sin is lawlessness.
					\end{description}
					
				1 John 3:4 has two parts.
				In the first part it tells us something about a person who sins;
					in the second part,
						it gives a somewhat abstract description of sin.
				The key word in both parts is \gr{ἀνομία}.
				In its form the word literally meanings ``lawlessness,''
					as it is a combination of the Greek privative prefix \gr{ἀ}- and the nominal ending -\gr{νομία},
						which itself derives from \gr{νόμος},
							``law.''
				In English we could say that \gr{ἀνομία} means ``not-lawfulness,''
					but the most obvious translation would be ``lawlessness.''
					
				The KJV renders the word as a phrase,
					``transgression of the law.''
				This translation shows that laws are involved,
					but with the English ``lawlessness'' available,
						there's no reason to introduce the phrase.
				The NASB,
					NIV,
						and NRSV also just have ``lawlessnes.''
				(The Anchor Bible volume discusses the issue of ``lawlessness'' and ``iniquity'' at 398--400.)
				
				On this reading,
					the second part of the verse is a description of sin that almost amounts to a definition:
						\gr{ἡ ἁμαρτία ἐστὶν ἡ ἀνομία},
							``sin is lawlessness.''
				For the whole thing,
					I like the NASB translation with present participles,
						to emphasize the continuous aspect:
				
					\begin{quote}
						Everyone who is practicing sin is also practicing lawlessness; and sin is lawlessness.
					\end{quote}
					
				This explicit link between sin and law is very useful.
				For one,
					it allows us to make better justified connections to important passages in restoration scripture,
						such as \hyperref[ALMA-42-17]{Alma 42:17} (``Now how could a man repent except he should sin? How could he sin if there was no law?'') and \hyperref[DC88-35]{DC 88:35} (``That which breaketh a law, and abideth not by law, but seeketh to become a law unto itself, and willeth to abide in sin, and altogether abideth in sin, cannot be sanctified by law, neither by mercy, justice, nor judgment'').
				I will discuss these connections later.
				
				Another benefit of the link is that it helps us better see how to understand righteousness.
				Later I will discuss the relation between sin and righteousness as well,
					but for now it's enough to say that:
						given sin as lawlessness,
							and given that righteousness is the opposite of sin,
								we should therefore see righteousness as ``lawfulness.''
				In fact,
					both BDAG and LSJ note that \gr{δικαιοσύνη},
						a common word in the New Testament almost always translated as ``righteousness,''
							is the proper opposite of \gr{ἀνομία}.
				Righteousness is lawful conduct;
					sin is unlawful conduct.
				Some verses in the New Testament make this relation clear,
					such as \hyperref[2CORINTHIANS-6-14]{2 Corinthians 6:14},
						part of which asks,
							\gr{τίς γὰρ μετοχὴ δικαιοσύνῃ καὶ ἀνομίᾳ}?
				``What do righteousness and lawlessness have in common?''
				More on that later.
				
			%
			\subparagraph{\hyperref[1JOHN-5-17]{1 John 5:17}}
			%
				\label{SIN-1john-5-17}
				
				The third verse to look at is 1 John 5:17.
				This one will be similar to the previous verse we saw in 1 John.
				Here is the KJV and Greek:
				
					\begin{description}
						\item[KJV] All unrighteousness is sin: and there is a sin not unto death.
						\item[SBLG] \gr{πᾶσα ἀδικία ἁμαρτία ἐστίν, καὶ ἔστιν ἁμαρτία οὐ πρὸς θάνατον.}
					\end{description}
					
				\noindent The part we care about is at the beginning,
					translated in English as ``all unrighteousness is sin.''
				Here are some other translations:
				
					\begin{description}
						\item[Anchor Bible] All wrongdoing is sin, but not all sin is deadly.
						\item[Syriac-English] For every wrongdoing is sin, but not every sin leads to death.
						\item[NASB] All unrighteousness is sin, and there is a sin not leading to death.
						\item[NIV] All wrongdoing is sin, and there is sin that does not lead to death.
						\item[NRSV] All wrongdoing is sin, but there is sin that is not mortal.
					\end{description}
					
				\noindent As before,
					the key word here is the one connected to \gr{ἁμαρτία} via the linking verb \gr{ἐστίν}.
				That word is \gr{ἀδικία}.
				The KJV and NASB render it as ``unrighteousness,''
					while the other three translations from the Greek give it as ``wrongdoing.''
					
				These are fine translations.
				Looking at the word in more detail,
					though,
						will give us more to go on in understanding sin.
				
				According to BDAG,
					\gr{ἀδικία} can be either an action ``that violates standards or right conduct,''
						or ``the quality of injustice'' itself,
							which would be a more abstract meaning with less reference to individual actions.
				BDAG also notes,
					however,
						that \gr{ἀδικία} is another opposite of \gr{δικαιοσύνη},
							``righteousness.''
				In fact it's similar to \gr{ἀνομία} that we saw before in that it begins with a privative prefix,
					\gr{ἀ}-,
						followed by -\gr{δικία},
							a nominal ending related to \gr{δίκαιος} and \gr{δικαιοσύνη}.
				We know from before that \gr{δικαιοσύνη} is lawful conduct,
					and the opposite of that would be a wrong that is unlawful.
				So with \gr{ἀδικία} in John 5:17,
					it isn't hard to make the connection to laws and our obedience to them,
						as we did with John 3:4.
				
				According to John 5:17,
					sin is an action that falls outside of whatever the conditions of \gr{δικαιοσύνη} require---%
						those conditions involve doing good,
							always within the laws set for us,
								so that we don't practice \gr{ἀνομία}.
								
		% --
		\paragraph{Sin, law, and righteousness}
		% --
		
			We've now covered three verses which help us understand how to think about sin.
			While it may be that none is a true definition of sin,
				each teaches us something about it.
			We can combine what we learn from these verses with other things the scriptures teach us to flesh out our understanding.
			
			The first connection to deepen is that between sin and law.
			Two of the scriptures we looked at dealt with law:
				1 John 3:4 told us that sin is lawlessness,
					and 1 John 5:17 told us that all \gr{ἀδικία} is sin,
						another word connected to laws.
			
			Many other scriptures discuss the connection between sin and law.
			On the one hand we should expect such detailed discussions in the scriptures,
				because in much of the Old Testament,
					for example,
						to sin was just to violate the law of Moses.
			Taken in the context of that law,
				it wouldn't be surprising that some people considered sin to be lawlessness.
				
			Many passages in \hyperref[ROMANS]{Romans} describe the connection this way.
			\hyperref[ROMANS-5-13]{Romans 5:13} says that ``sin is not imputed when there is no law,''
				meaning that sin isn't charged to someone's account---%
					it doesn't count as sin for that person---%
						when there isn't a law.
			\hyperref[ALMA-42-17]{Alma 42:17} makes a similar point,
				asking ``How could he [any person] sin if there was no law?''
			When a person knows the law,
				their knowledge makes them aware of sin,
					thus making them accountable for what they know---%
						``for by the law is the knowledge of sin'' (\hyperref[ROMANS-3-20]{Romans 3:20}),
							and ``as many as have sinned in the law shall be judged by the law'' (\hyperref[ROMANS-2-12]{Romans 2:12}).
			Sin is lawlessness,
				but ``lawless'' actions only count as sins if we're aware of the law.
			Christ's sacrifice covers those actions done by people ``who have ignorantly sinned'' (\hyperref[MOSIAH-3-11]{Mosiah 3:11}).
			
			The fact that the Book of Mormon speaks on this topic should encourage us to understand the ``laws'' as going beyond the law of Moses.
			The connections between sin and law get more enlightening as we do.
			
			The connection I want to make is between sin,
				as understood by lawlessness or lawless conduct,
					and Brigham Young's understanding of the priesthood.
			It has become common for people to explain the priesthood by saying something like,
				``It is God's power, given to man, for the benefit of His children.''
			There's nothing wrong with this,
				but Brigham Young on many occasions outlined a much more interesting view of what the priesthood is.
			For him,
				the priesthood is the system of government,
					or system of laws,
						which governs not only individuals and groups of people but also matter and physical interactions.
			Here are some representative quotations:
			
				\begin{description}
					\item[{\hyperref[EMS-1854-BY-12-FEB-1854]{12 February 1854}}] I recolect once when I was preaching a question was asked me, ``What is \sout{the} Preisthood.'' The answer was ready, and perfectly simple in its nature, \sout{and} plain to be understood, <and> \sout{\&} \sout{is} couched in a a short sentence (viz) Preisthood is a perfect system of Government that rules and reigns in eternity...The eternal preisthood of God,---the Gov\supers{t.} of God,---the laws of eternity, <is> \sout{are} a pure and perfect system of Gov\supers{t.}
					\item[{\hyperref[EMS-1854-BY-27-JUN-1854]{27 June 1854}}] The Priesthood is a pure system of government which governs the world and a small portion has come down to earth. It is the constitutional laws and pure system by which men are to be governed. 
				\end{description}
			
			don't just have to think of this in terms of the law of moses; in fact it's more powerful if you go beyond that
			
			See also: \hyperref[MATTHEW-5-20]{Matthew 5:20} (look at the Greek here too); 
				% BY quotes on priesthood and law
			
			
		% "Power" to forgive sins
			% Matthew 9:5-6; Mark 2:10; Luke 5:24. In all three verses the "power" is ἐξουσία.
			% Compare Helaman 5:10--1
		% Being a "servant" to sin or in "bondage"
			% John 8:34; the key part is πᾶς ὁ ποιῶν τὴν ἁμαρτίαν δοῦλός ἐστιν τῆς ἁμαρτίας·
			% Romans 6:16--20; similar language with forms of δοῦλός
			% Romans 7:18-25; 8:15; 8:21
			% 1 John 3:8, 10
			% dc 84:49--53
			% DC 138:57
		%XXX % Being "without law" (Romans 2:12) has important consequences for whether actions are sins; sin is intimately related the law and one's awareness of it
			%XXX % Romans 3:20, the last part is διὰ γὰρ νόμου ἐπίγνωσις ἁμαρτίας.
			%XXX % See also Romans 5:13
		% The whole "count/account/impute" thing, in Romans 4 and elsewhere; I'm not sure what to do with this yet.
			% Where there is no law, there is no transgression; see Romans 5:13.
			% alma 42:16--24, very important here
				% for ``repentance granted,'' think conditions of repentance; see my study in ``condition''
				% there's just one big set of laws, but we can refer to some of them as ``mercy'' because they have merciful outcomes, others as ``justice,'' and so on\
				% see also moroni 8:25--26
			% DC 88:32--40 also pairs well with Alma 42
				% see the note at DC 88:34 in particular
			% also DC 121
			% 2 nephi 2:13, 23---for 23, they were not yet hooked into the system of laws
			% the conscious awareness of what the law is counts for a lot---see mosiah 3:11---there is something fundamental about knowing the law for yourself
				% moroni 8:11
				% dc 29:46--47
					% dc 68:25
				% compare the thing about knowing the will of god vs. mere belief
					% 3 nephi 6:18
		% 1 corinthians 15:56, the sting of death/the strength of sin
			% Death "stings" (has power) because of sin, but sin only has power because of the law---this is a broad principle.
			% The word is δύναμις---the effectiveness, power, efficacy, something like that, of sin is that the law exists. Alma 42:22 here also.
		% Death is the result of sin---that's what is decreed in the laws
			% James 1:15
		% Definitions of sin
			% James 4:17
			% 1 John 3:4
			% 1 John 5:17
			% (don't forget about 1 corinthians 15:56 here)
		% alma 41:10, wickedness never was happiness---the sets of laws are different. People are malleable enough that you can get happiness in wickedness for a while but not in the long run.
			% Mormon 2:13
		
		
		
		
		
		
		
		
		
		
		
		
		
		
	% -----
	\subsubsection{The difference between sin and transgression}
	% -----
		\label{SIN-sin-transgression}
		
		There are many places in the BOM where it appears as though the terms ``sin'' and ``transgression'' are interchangeable. \hyperref[JAROM-1-10]{Jarom 1:10} is an example, and \hyperref[ALMA-9-14]{Alma 9:14} is another.
		
		But there are places where they do not seem substitutable for each other. The most prominent example is when the BOM and other scriptures speak of Adam's transgression---the scriptures \textit{never}, as far as I can tell in my searching, refer to the events surrounding the fall as ``sin'' or ``a sin.'' They \textit{always} call it a ``transgression.'' The consistency is striking, and it's maintained even in cases where the terms ``sin'' and ``transgression'' appear in close proximity, such as \hyperref[ROMANS-5-14]{Romans 5:14} and \hyperref[MOSIAH-3-11]{Mosiah 3:11}.
		
		And of course, there are clear differences in the Greek words behind some of these passages; sin is \hyperref[HAMARTIA]{\gr{ἁμαρτία}}, while transgression is \hyperref[PARABASIS]{\gr{παράβασις}}.
		
		I think the key to discerning the difference between sin and transgression is in their relation to laws. Transgressions have to do more directly with laws and breaking laws. As we read in the \hyperref[TRANSGRESSION-BD-TDNTA]{TDNTA}, there is transgression only where there is law; the presence of a law is a necessary condition for a transgression to be possible. When there is no law, there is no transgression.
		
		But if you think of sin as I defined it above---as being able to do good, but not doing it---and if you think of the consequences of sin as a decreased measure of the Spirit, power, and knowledge, then it's possible to see sin as having less to do with specific laws and more to do with how we use our agency to develop ourselves. Sin is a more open-ended notion, in a way; the choice to sin is the choice to forego some learning or growth, whether or not a specific law is involved. 
		
		Maybe I can put it this way, although I need to think about this one a bit more: when you transgress, the law condemns you; when you sin, you condemn yourself. Sinning is saying, ``No more growth for me'' (or ``No growth right now, please,'') whatever the ``law'' might say.
		
		This way of seeing the difference between sin and transgression casts \hyperref[DC58-26]{DC 58:26--29} into a bit of a new light:
		
			\begin{quote}
				For behold, it is not meet that I should command in all things; for he that is compelled in all things, the same is a slothful and not a wise servant; wherefore he receiveth no reward.
				27 Verily I say, men should be anxiously engaged in a good cause, and do many things of their own free will, and bring to pass much righteousness;
				28 For the power is in them, wherein they are agents unto themselves. And inasmuch as men do good they shall in nowise lose their reward.
				29 But he that doeth not anything until he is commanded, and receiveth a commandment with doubtful heart, and keepeth it with slothfulness, the same is damned.
			\end{quote}
			
		We can take ``It is not meet that I should command in all things'' to be something like, ``I shouldn't have to give a law for everything.'' This is because we should be ``anxiously engaged'' in actually \textit{doing good}---because not doing that would be sin, according to the definition. We have enough power, or enough capacity, to bring about goodness and righteousness without being told exactly how to do that. On the negative side, if we have to wait until we have a law, then we're damned, in the sense that we are choosing not to do good when we could do it, and are therefore sinning. We have condemned ourselves by forgoing the power, knowledge, and growth that we could have gained by doing the good ourselves.
		
		These verses encourage us to think in terms of, what good can I do? This is a growth mindset, one that embraces the opportunity to learn and to change. The verses discourage us from thinking in terms like, what have I been commanded to do? That's a mindset whose goal is just to avoid transgression, but that isn't good enough. We need to be growing and expanding beyond the particular commandments that we have been given. 
		
		Note, also, that sin and transgression won't come apart like this all the time. As I said above, there are places in the scriptures where they are effectively the same thing, meaning that in many situations in our lives they are the same. But they are different concepts and in some situations their differences are revealed.
		
		As some things to follow up on in the future, consider relating transgression with a remission of sins---being free from the consequences of broken law---and also therefore with justification. You might also relate sin, or the growth mindset of the sin-game, with sanctification (maybe, just to see what comes out).